\section*{अध्याय \ref{ch:b2:reaction-free-energy} --- अभिक्रिया की एन्ट्रॉपी और मुक्त ऊर्जा}

\begin{solution}{exo:b2:reaction-free-energy:1}
(a) गैस के तीन मोल द्रव बनते हैं: $\Delta_r S^\circ < 0$। (b) एक गैस
बनती है: $> 0$। (c) गैस का एक अणु दो देता है: $> 0$। (d) विलयन के आयन
एक क्रिस्टल बनाते हैं: $< 0$। (e) गलन: $> 0$। (f) $\Delta\nu_{\text{गैस}} = 0$:
छोटा, और यहाँ धनात्मक (सारणियों से \qty{+20}{J/(K.mol)}, क्योंकि \ce{HCl}
\ce{H2} और \ce{Cl2} से कम सममित है)।
% ledger: s0:HCl_g, s0:H2_g, s0:Cl2_g
\end{solution}

\begin{solution}{exo:b2:reaction-free-energy:2}
$\Delta_r S^\circ = 2(192.77) - 191.61 - 3(130.68) = \qty{-198.1}{J/(K.mol)}$;
$\Delta_r G^\circ = -91.9 - 298.15 \times (-0.1981) = \qty{-32.8}{kJ/mol}$:
\qty{298}{K} पर \omterm{def:b2:reaction-free-energy:exergonic}{ऊर्जामोची}, यद्यपि एन्ट्रॉपी पद (\qty{+59.1}{kJ/mol})
इसके विरुद्ध काम करता है।
\end{solution}

\begin{solution}{exo:b2:reaction-free-energy:3}
$\Delta_r S^\circ = 69.95 - 130.68 - \tfrac12(205.15) = \qty{-163.3}{J/(K.mol)}$;
$\Delta_r G^\circ = -285.83 + 298.15 \times 0.1633 = \qty{-237.14}{kJ/mol}$।
यह द्रव जल की संभवन अभिक्रिया है (संदर्भ अवस्थाओं वाले तत्वों से एक मोल),
इसलिए इसकी \omterm{def:b2:reaction-free-energy:reaction-gibbs}{मानक अभिक्रिया गिब्स ऊर्जा} परिभाषा से ही
$\Delta_f G^\circ(\ce{H2O}, \text{l})$ है; सारणियाँ भी यही
मान देती हैं।
\end{solution}

\begin{solution}{exo:b2:reaction-free-energy:4}
\qty{298}{K} पर लगभग \qty{42}{J/(K.mol)} और \qty{1000}{K} पर
\qty{87}{J/(K.mol)} (द्रव)। उछाल: \qty{692.7}{K} पर \qty{10.6}{J/(K.mol)} और
\qty{1180.2}{K} पर \qty{97.7}{J/(K.mol)}। $\Delta_{\text{गलन}}H = 10.6 \times
692.7 = \qty{7.3}{kJ/mol}$ और $\Delta_{\text{वाष्पन}}H = 97.7 \times 1180.2 =
\qty{115}{kJ/mol}$।
\end{solution}

\begin{solution}{exo:b2:reaction-free-energy:5}
$\Delta_r H^\circ = \qty{44.00}{kJ/mol}$, $\Delta_r S^\circ = 188.84 - 69.95 =
\qty{118.9}{J/(K.mol)}$, $\Delta_r G^\circ = 44.00 - 298.15 \times 0.1189 =
\qty{+8.55}{kJ/mol}$। \qty{298}{K} पर परिवर्तन उत्क्रमणीय नहीं है (द्रव
जल और \qty{1}{bar} पर वाष्प साम्य में नहीं हैं), इसलिए $\Delta S \ne
Q/T$: $\Delta_r H^\circ/T = \qty{147.6}{J/(K.mol)}$। धनात्मक
$\Delta_r G^\circ$ का अर्थ है कि \qty{298}{K} पर जल \qty{1}{bar} पर शुद्ध वाष्प के
वायुमंडल में वाष्पित नहीं होता: वाष्प संघनित हो जाती है। दोनों
वहाँ साम्य में होते हैं जहाँ $\Delta_r G^\circ = 0$: $T = 44\,004/118.9 =
\qty{370}{K}$, सामान्य क्वथनांक (\qty{373}{K}) से \qty{3}{K} के भीतर:
\qty{75}{K} के अंतराल में \omterm{def:b2:reaction-free-energy:ellingham-approximation}{एलिंगम सन्निकटन} अच्छा है।
\end{solution}

\begin{solution}{exo:b2:reaction-free-energy:6}
$T_i = 57\,100/175.7 = \qty{325}{K}$। इससे नीचे $\Delta_r G^\circ > 0$ और
रंगहीन \ce{N2O4} अनुकूलित होता है; नली को ठंडा करने से गैस
\ce{N2O4} की ओर विस्थापित होती है और \ce{NO2} का भूरा रंग फीका पड़ता है (संघटन से
मात्रात्मक संबंध \cref{ch:b2:equilibrium-shifts} में है)।
\end{solution}

\begin{solution}{exo:b2:reaction-free-energy:7}\emergencystretch=3em
$\Delta_r S^\circ = -\dd\Delta_r G^\circ/\dd T = -(-4.0 + 12.0)/100 =
\qty{-0.080}{kJ/(K.mol)} = \qty{-80}{J/(K.mol)}$ और $\Delta_r H^\circ =
\Delta_r G^\circ + T\Delta_r S^\circ = -12.0 + 300 \times (-0.080) =
\qty{-36.0}{kJ/mol}$। जाँच: $\Delta_r G^\circ/T$, $-0.0400$ से
$\qty{-0.0100}{kJ/(K.mol)}$ तक जाता है, प्रवणता प्रति केल्विन $3.00 \times 10^{-4}$, जो
$-\Delta_r H^\circ/T^2 = 36.0/346^2 = 3.00 \times 10^{-4}$ के बराबर है, ज्यामितीय
माध्य ताप $\sqrt{300 \times 400} = \qty{346}{K}$ पर।
\end{solution}

\begin{solution}{exo:b2:reaction-free-energy:8}\emergencystretch=3em
$\Delta_r S^\circ = 186.25 - 5.74 - 2(130.68) = \qty{-80.8}{J/(K.mol)}$,
$\Delta_r G^\circ = -74.87 + 298.15 \times 0.0808 = \qty{-50.8}{kJ/mol}$,
जो मेथेन का सारणीबद्ध $\Delta_f G^\circ$ है। दोनों पदों का चिह्न
एक ही होने से, मेथेन $T_i = 74\,870/80.8 \approx
\qty{926}{K}$ से ऊपर अस्थायी हो जाती है (\omterm{def:b2:reaction-free-energy:ellingham-approximation}{एलिंगम सन्निकटन}): गर्म मेथेन कार्बन और
हाइड्रोजन में भंजित होती है, जो दोनों को पाने का एक मार्ग है।
\end{solution}

\begin{solution}{exo:b2:reaction-free-energy:9}
प्रतिस्थापन: \ce{CH3CHBrCH3 + OH- -> CH3CH(OH)CH3 + Br-}; विलोपन:
\ce{CH3CHBrCH3 + OH- -> CH2=CHCH3 + H2O + Br-}। दोनों मानक अभिक्रिया
गिब्स ऊर्जाओं का अंतर $\Delta(\Delta_r G^\circ) = 40 - T \times
0.090$ है, जो $40\,000/90 \approx \qty{440}{K}$ से ऊपर ऋणात्मक है। विलोपन
प्रतिस्थापन से एक अणु अधिक बनाता है: उसकी एन्ट्रॉपी अधिक है, और एन्ट्रॉपी
पद $T$ के अनुपात में बढ़ता है।
\end{solution}

\begin{solution}{exo:b2:reaction-free-energy:10}
$N_A$ अणुओं में से हर एक के 2 अभिविन्यास हैं: $W = 2^{N_A}$, अतः $S =
k_B\ln 2^{N_A} = R\ln 2 = \qty{5.76}{J/(K.mol)}$। क्रिस्टल
पूर्णतः व्यवस्थित नहीं है: निम्न ताप पर अणु अपने यादृच्छिक अभिविन्यासों में
जम जाते हैं और उस एकमात्र व्यवस्थित विन्यास तक नहीं पहुँच पाते जिसे तृतीय
नियम मानता है।
\end{solution}

\begin{solution}{exo:b2:reaction-free-energy:11}
$\dd\Delta_r S^\circ/\dd T = c/T$ का समाकलन कीजिए: $\Delta_r S^\circ(T) = \Delta_r
S^\circ(T_0) + c\ln(T/T_0)$; किरखॉफ़ से $\Delta_r H^\circ(T) = \Delta_r
H^\circ(T_0) + c(T - T_0)$; $T\Delta_r S^\circ(T)$ घटाइए। चूना पत्थर के लिए
\qty{1100}{K} पर: एलिंगम $178.32 - 1100 \times 0.16064 = \qty{1.62}{kJ/mol}$;
संशोधन $c[T - T_0 - T\ln(T/T_0)] = -0.00195 \times (801.85 - 1436.3) =
\qty{+1.24}{kJ/mol}$: $\Delta_r G^\circ = \qty{2.86}{kJ/mol}$। त्रुटि,
\qty{1.2}{kJ/mol}, $\Delta_r H^\circ$ के एक प्रतिशत से कम है, परंतु
\omterm{def:b2:reaction-free-energy:inversion-temperature}{व्युत्क्रमण ताप} के निकट यही चिह्न तय करती है।
\end{solution}

\begin{solution}{exo:b2:reaction-free-energy:12}
(a) $\Delta_r H^\circ = -763.2 + 944.0 = \qty{+180.8}{kJ/mol}$, $\Delta_r
S^\circ = 353.2 + 205.15 - 50.62 - 2(223.08) = \qty{+61.6}{J/(K.mol)}$;
$\Delta_r G^\circ(1000) = 180.8 - 61.6 = \qty{+119.2}{kJ/mol}$: \omterm{def:b2:reaction-free-energy:exergonic}{ऊर्जाग्राही},
क्लोरीनन नहीं होता। (b) \ce{C + O2 -> CO2} के लिए $\Delta_r G^\circ(1000) =
-393.5 - 2.9 = \qty{-396.4}{kJ/mol}$। योग, \ce{TiO2(s) + 2Cl2(g) + C(s)
-> TiCl4(g) + CO2(g)}, के लिए $\Delta_r G^\circ = \qty{-277.2}{kJ/mol}$: कार्बन
पहली अभिक्रिया द्वारा बनी डाइऑक्सीजन का उपभोग कर लेता है, और युग्मित अभिक्रिया
प्रबल रूप से \omterm{def:b2:reaction-free-energy:exergonic}{ऊर्जामोची} है।
\end{solution}

\begin{solution}{pb:b2:reaction-free-energy:1}
\textbf{1.} \ce{CaCO3(s) -> CaO(s) + CO2(g)}।
\textbf{2.} $\Delta_r H^\circ = -635.09 - 393.51 + 1206.92 =
\qty{+178.3}{kJ/mol}$: \omterm{def:b2:reaction-enthalpy:exothermic}{ऊष्माशोषी}।
\textbf{3.} $\Delta_r S^\circ = 39.75 + 213.79 - 92.9 = \qty{+160.6}{J/(K.mol)}$,
धनात्मक, क्योंकि ठोस से एक गैस बनती है।
\textbf{4.} $\Delta_r G^\circ = 178.32 - 298.15 \times 0.16064 =
\qty{+130.4}{kJ/mol}$।
\textbf{5.} हाँ: विघटन के लिए $\Delta_r G < 0$ आवश्यक है, और यहाँ यह
प्रबल रूप से धनात्मक है; वायु की कार्बन डाइऑक्साइड, जो \qty{1}{bar} से बहुत नीचे है,
$\Delta_r G$ को घटाती तो है, पर \qty{130}{kJ/mol} से कहीं कम (अगला अध्याय)।
\textbf{6.} $\Delta_r C^\circ_p = 42.80 + 37.13 - 81.88 =
\qty{-1.95}{J/(K.mol)}$, बहुत छोटा: $\Delta_r H^\circ$ और $\Delta_r S^\circ$
$T$ के साथ मुश्किल से बदलते हैं।
\textbf{7.} $\Delta_r G^\circ(T) = 178.32 - 0.16064\,T$ (\unit{kJ/mol})।
\textbf{8.} \qty{1000}{K} पर \qty{+17.7}{kJ/mol}; \qty{1300}{K} पर
\qty{-30.5}{kJ/mol}।
\textbf{9.} $T_i = 178\,320/160.64 = \qty{1110}{K}$।
\textbf{10.} $\Delta_r G^\circ/T = 178.32/T - 0.16064$, जिसका अवकलज
$-178.32/T^2 = -\Delta_r H^\circ/T^2$ है।
\textbf{11.} $T_i$ पर, $c[T_i - T_0 - T_i\ln(T_i/T_0)] = -0.00195 \times
(811.9 - 1459.4) = \qty{+1.26}{kJ/mol}$; $\Delta_r G^\circ$ शून्य तक
$1260/160.6 \approx \qty{8}{K}$ ऊपर, लगभग \qty{1118}{K} पर, पहुँचता है: एक
प्रतिशत से कम।
\textbf{12.} \qty{1000}{K} पर $\Delta_r G > 0$: विघटन आगे नहीं
बढ़ सकता (चूना कार्बन डाइऑक्साइड अवशोषित करेगा)। \qty{1300}{K} पर $\Delta_r G <
0$: चूना पत्थर विघटित होता है।
\textbf{13.} $\Delta_r G$ धनात्मक हो जाता है: चूना कार्बन डाइऑक्साइड ग्रहण करता है
और फिर से कार्बोनेट बन जाता है।
\textbf{14.} $\delta S_{\text{c}} = -\Delta_r G\,\dd\xi/T = 30\,506/1300 =
\qty{23.5}{J/K}$ प्रति मोल।
\textbf{15.} हाँ: \cref{ch:b2:chemical-potential} दिखाएगा कि
$\Delta_r G = \Delta_r G^\circ + RT\ln(p_{\ce{CO2}}/p^\circ)$, जो
कार्बन डाइऑक्साइड का दाब \qty{1}{bar} से नीचे होने पर घटता है, जिससे
विघटन $T_i$ से नीचे ही आरंभ हो जाता है।
\textbf{16.} $n = 10^6/56.08 = \qty{1.783e4}{mol}$।
\textbf{17.} $1.783 \times 10^4 \times 178.32 = \qty{3.18e6}{kJ}$, अर्थात्
प्रति टन \qty{3.18}{GJ}।
\textbf{18.} $1.783 \times 10^4 \times 44.01 = \qty{785}{kg}$ \ce{CO2}।
\textbf{19.} दी जाने वाली ऊष्मा $3.18 \times 10^6/0.50 = \qty{6.36e6}{kJ}$, अतः
$6.36 \times 10^6/802.3 = \qty{7.93e3}{mol}$ मेथेन, \qty{127}{kg}।
\textbf{20.} $7.93 \times 10^3 \times 44.01 = \qty{349}{kg}$ \ce{CO2}।
\textbf{21.} $785 + 349 = \qty{1134}{kg}$ प्रति टन चूना, जिसका 69\,\% चूना पत्थर
से आता है।
\textbf{22.} चूना पत्थर की कार्बन डाइऑक्साइड रससमीकरणमिति से तय होती है,
भट्ठी को चाहे जो गर्म करे; केवल ईंधन का अंश घटाया जा सकता है।
\textbf{23.} चूना पत्थर \qty{1}{bar} कार्बन डाइऑक्साइड के अंतर्गत अपने
\omterm{def:b2:reaction-free-energy:inversion-temperature}{व्युत्क्रमण ताप}, $\boldsymbol{T_i \approx \qty{1.11e3}{K}}$
(लगभग \qty{840}{\celsius}), से ऊपर विघटित होता है।
\end{solution}
